\subsection{Application to composite extensions}

PROPOSITION 10. -- Let L and M be two extensions of a field K and (E, u, v) a composite extension of L and M (V, p.~12). Suppose that the ring L ⊗\_K M is an integral domain and that the subextensions u(L) and v(M) of E are algebraically disjoint over K. Then u(L) and v(M) are linearly disjoint over K.

Put \(w = u * v\) (V, p.~12), denote by F the field of fractions of the integral domain \(L \otimes_{K} M\) and identify L (resp. M) with a subfield of F by means of the mapping \(x \mapsto x \otimes 1\) (resp. \(y \mapsto 1 \otimes y\) ); then the restriction of w to L (resp. M) is u (resp. v). Let B be a transcendence basis of M over K (V, p.~109, Th. 1).

By hypothesis \(u(\mathbf{L})\) and \(v(\mathbf{M})\) are algebraically disjoint over \(\mathbf{K}\) ; therefore (V, p.~114, Prop. 14), \(u(\mathbf{L})\) and \(v(\mathbf{K}(\mathbf{B}))\) are linearly disjoint over \(\mathbf{K}\) . Thus there exists a K-homomorphism \(u': \mathbf{L}(\mathbf{B}) \to \mathbf{E}\) which agrees with \(u\) on \(\mathbf{L}\) and with \(v\) ou \(\mathbf{K}(\mathbf{B})\) . By construction \(\mathbf{L}\) and \(\mathbf{M}\) are linearly disjoint over \(\mathbf{K}\) in \(\mathbf{F}\) ; by Prop. 8 (V, p.~15) the subfields \(\mathbf{L}(\mathbf{B})\) and \(\mathbf{M}\) of \(\mathbf{F}\) are linearly disjoint over \(\mathbf{K}(\mathbf{B})\) . It follows that there exists a K-homomorphism \(w' \colon \mathbf{M}[\mathrm{L}(\mathrm{B})] \to \mathrm{E}\) which agrees with \(u'\) on \(\mathrm{L}(\mathrm{B})\) and with \(v\) on \(\mathbf{M}\) . But the field \(F\) is generated by \(\mathbf{M} \cup \mathrm{L}(\mathrm{B})\) and \(\mathbf{M}\) is algebraic over \(\mathbf{K}(\mathbf{B})\) ; we thus have \(\mathbf{M}[\mathrm{L}(\mathbf{B})] = F(\mathbf{V}, p. 18, Cor. 2)\) . Hence we conclude that \(w'\) is a K-isomorphism of \(F\) onto \(E\) whose restriction to \(L\) (resp. \(M\) ) is \(u\) (resp. \(v\) ). This shows \(u(L)\) and \(v(M)\) to be linearly disjoint over \(K\) .

COROLLARY 1. --- Let \(\Omega\) be an extension of a field \(K\) and \(L\) a subextension of \(\Omega\) which is regular over \(K\) . Every subextension \(M\) of \(\Omega\) which is algebraically disjoint from \(L\) over \(K\) is linearly disjoint.

The ring \(L \otimes_{K} M\) is an integral domain by definition of regular extension and it suffices now to apply Prop. 10.

COROLLARY 2. --- Let \(\Omega\) be an extension of a field \(K\) and \(L\) , \(M\) two subextensions of \(\Omega\) . Suppose that \(L\) is separable over \(K\) and that the relative separable closure of \(K\) in \(M\) is equal to \(K\) . If \(L\) and \(M\) are algebraically disjoint over \(K\) , then they are linearly disjoint over \(K\) .

By Prop. 10 it is enough to remark that the ring \(\mathbf{L} \otimes_{\mathbf{K}} \mathbf{M}\) is an integral domain (V, p.~140, Cor.).

